\section{问题重述}

\subsection{问题背景}

研究对象是一个经公共连接点接入外网的小区级能源系统。光伏出力具有间歇性，居民用电具有日周期，电池则把不同时刻的能量联系起来；三者共同决定每个 $10$ min 时段应从外网买入多少电。调度既要逐时满足用电需求，又要利用电池在低价或光伏富余时吸收能量、在缺电或高价时释放能量，因而本质上是带状态递推、预测误差和偏差结算的动态优化问题。

本问题的物理对象由三部分组成：光伏发电单元、储能设备与小区负载，并通过公共连接点与外网相连。附录 1 给定储能参数：最大容量 $12000$ kWh，最大充放电功率 $5000$ kW，2025 年 1 月 1 日 0:00 储电量为 $6000$ kWh，运行期储电量必须保持在 $1200\sim10800$ kWh 之间，充放电效率均为 $90\%$。附件数据的时间分辨率均为 $10$ min，即一天可分为 $144$ 个调度时段。

问题的核心困难在于\textbf{信息的可获得性}：微网必须在决策时刻仅依据当时已知的信息制定策略，而负载、光伏与实际电价在一天内是逐步揭晓的。因此四个问题的差异并不在物理约束，而在于\emph{决策时刻能看到什么信息}，以及\emph{计划与实际不一致时如何结算}。

\subsection{问题提出}

\indent\textbf{问题一：}在典型日电价、负载和光伏预测均给定时，制定 $0{:}00$ 日前计划，并满足首末储电量相同。论文须按题面表 1、表 2 给出六个指定时段、全天购电量与费用，以及六个 4 h 分块的充放电量；完整结果写入 \texttt{result1.xlsx}。

\indent\textbf{问题二：}在负载和光伏逐日变化、计划缺口按 $5$ 倍电价紧急补购时，制定每日 $0{:}00$ 计划。论文须对 2025-03-20、06-21、09-23、12-21 按表 1--3 报告购电、储能和紧急购电结果；完整结果写入 \texttt{result2.xlsx}。

\indent\textbf{问题三：}每日四个发布时刻可获得未来 24 h 光伏预报，并允许调整 $0{:}00$ 基准购电量。需要写清计划费、上调费、下调违约费与紧急购电费，判断新增预报是否值得采用；同样对四个指定日期给出表 1--3，并写入 \texttt{result3.xlsx}。

\indent\textbf{问题四：}将附件 4 的逐时电价代入问题二、三重新求解，给出相应年度指标与四个指定日期结果，分别写入 \texttt{result4-2.xlsx} 和 \texttt{result4-3.xlsx}。

\subsection{研究现状综述}

微网经济调度问题在数学上通常被建模为以购电费或综合运行成本最小为目标的线性规划或混合整数规划，外网购电、分布式电源消纳与储能充放电被统一纳入同一组能量平衡约束中 \cite{wu2014dynamic,nemati2018optimization}。确定性模型适用于负载、电价与光伏均已知的日前场景，其优点是约束结构清晰、求解稳定可靠。

当光伏出力存在预测误差时，主流处理方式有三类：一是\textbf{情景法／两阶段随机规划}，用有限个历史或生成情景刻画不确定性，第一阶段确定与情景无关的日前计划，第二阶段按情景确定执行量 \cite{wu2016solution,sun2025distributionally}；二是\textbf{鲁棒优化}，在最坏情景下保证可行性，代价是结果偏保守 \cite{zhang2018robust}；三是\textbf{滚动时域优化（模型预测控制）}，在新信息到达后仅重优化剩余时域并只执行近期决策，天然适配“预报分批发布”的决策结构 \cite{elkazaz2020energy,silva2020optimal}。含光伏与储能的日前调度研究普遍指出，储能的价值高度依赖预测精度与调度时域长度 \cite{luo2020optimal}，而储能的运行成本又与其循环深度密切相关 \cite{maheshwari2020optimizing}。

针对“预测值应偏向哪一侧”的问题，报童模型给出了经典答案：当高估与低估的边际损失不对称时，最优的点预测不是条件均值，而是某个\textbf{分位数} \cite{qin2011newsvendor}。这一结论在电力系统中已被用于发电商的投标策略与偏差电量定价 \cite{polat2024renewable}。在预测方法层面，基于数值天气预报的机器学习模型是目前光伏功率预测的主流 \cite{markovics2022comparison}；在市场机制层面，实时电价与需求响应的耦合会显著改变用户的购电行为 \cite{wang2015review}。

现有研究多聚焦于算法本身，而对\textbf{信息集边界}与\textbf{结算口径}的刻画往往不够严格：在业务未提供全天计划负荷时，日前优化若直接使用当天真实负载，会混入理想信息优势；若已提供，则该情形可作为相应的信息口径。若比较策略时同时改变预测方法、约束和费用口径，差异也无法归因。本文据此把信息集边界作为建模原则，并用单因素消融与配对控制支持归因；求解则利用线性／混合整数线性结构，不依赖随机性较强的元启发式算法。
